\lyseis{Β3.2}{Μεσοκάθετος ευθύγραμμου τμήματος}
\sol{1} Ελεύθερη κατασκευή. Και με τους δύο τρόπους η ευθεία περνά από το σημείο που απέχει $3$ cm από το $\mathrm{Α}$ και σχηματίζει ορθές γωνίες με το $\mathrm{ΑΒ}$. Στον τρόπο με τον διαβήτη η ακτίνα πρέπει να είναι μεγαλύτερη από $3$ cm.
\sol{2} Ελεύθερη κατασκευή. Για κάθε σημείο της μεσοκαθέτου οι δύο αποστάσεις από τα $\mathrm{Α}$ και $\mathrm{Β}$ βγαίνουν ίσες (μέχρι τα σφάλματα της μέτρησης).
\sol{3} Κάθε σημείο της μεσοκαθέτου ισαπέχει από τα άκρα, άρα $\mathrm{ΜΒ}=\mathrm{ΜΑ}=5$ cm.
\sol{4} Πάνω στη μεσοκάθετο του $\mathrm{ΑΒ}$, γιατί κάθε σημείο που ισαπέχει από τα άκρα ανήκει σε αυτήν.
\sol{5} α) Σωστό. β) Λάθος: πρέπει να είναι και κάθετη στο τμήμα. γ) Σωστό. δ) Λάθος: ισαπέχουν μόνο τα σημεία της μεσοκαθέτου.
\sol{6} Ελεύθερη κατασκευή. Το $\mathrm{ΑΒ}$ χωρίζεται σε $4$ ίσα κομμάτια: το $\mathrm{Μ}$ το κόβει στη μέση και τα μέσα των $\mathrm{ΑΜ}$ και $\mathrm{ΜΒ}$ κόβουν το καθένα από τα δύο μισά στη μέση.
\sol{7} Ελεύθερη κατασκευή. Οι τρεις αποστάσεις βγαίνουν ίσες: το σημείο τομής ανήκει και στις δύο μεσοκαθέτους, άρα ισαπέχει από $\mathrm{Α}$, $\mathrm{Β}$ και από $\mathrm{Β}$, $\Gamma$, δηλαδή και από τα τρία σημεία.
\sol{8} Ελεύθερη κατασκευή. Γράφουμε δύο κύκλους ακτίνας $5$ cm με κέντρα $\mathrm{Α}$ και $\mathrm{Β}$. Τέμνονται σε δύο σημεία, που είναι οι δύο θέσεις. Βρίσκονται και τα δύο πάνω στη μεσοκάθετο του $\mathrm{ΑΒ}$, από τις δύο μεριές του $\mathrm{ΑΒ}$, σε ίση απόσταση από αυτό.
\sol{9} Φέρνουμε τη μεσοκάθετο του $\mathrm{ΑΒ}$. Το σημείο όπου τέμνει τον δρόμο $\varepsilon$ ισαπέχει από τα δύο χωριά, άρα εκεί ανοίγεται το πηγάδι. (Αν ο δρόμος $\varepsilon$ είναι παράλληλος στη μεσοκάθετο δεν υπάρχει τέτοιο σημείο.)
\sol{10} Φέρνουν τις μεσοκαθέτους των δύο ζευγαριών σπιτιών, π.χ. των $\mathrm{ΑΒ}$ και $\mathrm{ΒΓ}$. Το σημείο τομής τους ισαπέχει και από τα τρία σπίτια.
\sol{11} Φέρνει τις μεσοκαθέτους των $\mathrm{ΑΒ}$ και $\mathrm{ΒΓ}$. Το σημείο τομής τους ισαπέχει από τα τρία σημεία του χείλους, άρα είναι το κέντρο του καπακιού.
\sol{12} Πάνω στη μεσοκάθετο του τμήματος $\mathrm{ΑΒ}$ (σε οποιοδήποτε σημείο της).
\sol{13} Αν η ακτίνα ήταν μικρότερη από το μισό του $\mathrm{ΑΒ}$ οι δύο κύκλοι δεν θα έφταναν ο ένας τον άλλο και δεν θα είχαν κοινό σημείο. Αν ήταν ακριβώς το μισό, θα εφάπτονταν μόνο στο μέσο.
\sol{14} Η πτυχή είναι η μεσοκάθετος του $\mathrm{ΑΒ}$: αφού το $\mathrm{Α}$ πέφτει πάνω στο $\mathrm{Β}$, η πτυχή είναι άξονας συμμετρίας του τμήματος, άρα περνά από το μέσο του και είναι κάθετη σε αυτό.
\sol{15} Αν $\mathrm{ΑΒ}=6$ cm: δύο σημεία (οι δύο κύκλοι ακτίνας $5$ τέμνονται σε δύο σημεία, γιατί $5>3$). Αν $\mathrm{ΑΒ}=10$ cm: ένα σημείο, το μέσο του $\mathrm{ΑΒ}$ (οι κύκλοι εφάπτονται). Αν $\mathrm{ΑΒ}=12$ cm: κανένα (οι κύκλοι δεν τέμνονται).
\sol{16} Η $\mathrm{ΚΛ}$ είναι η μεσοκάθετος του $\mathrm{ΑΒ}$: περνά από το μέσο του και είναι κάθετη σε αυτό.
\sol{17} Ελεύθερη κατασκευή. Η πτυχή που προκύπτει πρέπει να συμπίπτει με τη μεσοκάθετο που σχεδίασες.
\sol{18} Ελεύθερη κατασκευή. Σχεδιάζουμε δύο χορδές $\mathrm{ΑΒ}$ και $\mathrm{ΒΓ}$ και τις μεσοκαθέτους τους. Το σημείο τομής είναι το κέντρο.
